\documentclass[11pt]{article}
\usepackage[margin=0.85in]{geometry}
\usepackage{lmodern}
\usepackage{microtype}
\usepackage{enumitem}
\usepackage{listings}
\usepackage[T1]{fontenc}

\lstset{
  basicstyle=\ttfamily\small,
  columns=fullflexible,
  keepspaces=true,
  showstringspaces=false,
  breaklines=true,
  frame=single,
  xleftmargin=0.4em,
  xrightmargin=0.4em,
  aboveskip=0.6em,
  belowskip=0.6em
}

\pagestyle{empty}
\setlength{\parskip}{0.4em}
\setlength{\parindent}{0pt}
\setlist[enumerate,1]{leftmargin=1.4em, itemsep=0.9em, topsep=0.4em}
\setlist[enumerate,2]{leftmargin=1.5em, itemsep=0.45em, topsep=0.3em, label=\arabic*.}
\setlist[itemize]{leftmargin=1.3em, itemsep=0.2em, topsep=0.25em}

\begin{document}

\begin{center}
{\LARGE\bfseries Midterm 2 --- Sample 2}\\[0.35em]
DATA 412/612
\end{center}

\vspace{0.3em}

Practice exam, questions only. Put your answers in an R Markdown document,
and make sure code is visible when you knit
(\texttt{knitr::opts\_chunk\$set(echo = TRUE)} in the first chunk).
Use the tidyverse rather than base R where you can.
You may use your notes, the R help files, and the internet.
Do not discuss the exam with anyone else.

This sample covers the same ground as Sample~1 (joins, strings, and
tidying), and it also uses factors and dates.

\begin{enumerate}
\item \textbf{Library visits} (10 pts).

\texttt{branches}: \texttt{branch} (name), \texttt{city}, \texttt{seats}.

\texttt{visits}: \texttt{patron}, \texttt{branch},
\texttt{check\_in} (date and time, stored as text),
\texttt{minutes} (length of the visit).

\begin{lstlisting}
library(tidyverse)

branches <- tribble(
  ~branch,    ~city, ~seats,
  "Shaw",     "DC",  80,
  "Petworth", "DC",  45,
  "Rosslyn",  "VA",  60
)

visits <- tribble(
  ~patron, ~branch,    ~check_in,             ~minutes,
  "Ari",   "Shaw",     "Mar 03, 2024 9:15am",  40,
  "Bea",   "Rosslyn",  "Mar 03, 2024 2:05pm",  15,
  "Cal",   "Petworth", "Mar 04, 2024 11:00am", 90,
  "Dee",   "Shaw",     "Mar 04, 2024 6:40pm",  25
)
\end{lstlisting}

\begin{enumerate}
\item (3 pts) Use \texttt{parse\_datetime()} to parse \texttt{check\_in},
and assign the result back to \texttt{visits}.

\item (3 pts) Add \texttt{city} and \texttt{seats} from \texttt{branches}
onto \texttt{visits}.

\item (4 pts) Add a column \texttt{weekday}: an ordered factor for the
weekday of \texttt{check\_in}, with levels Monday through Sunday in that
order.
\end{enumerate}

\item \textbf{Strings} (10 pts).
Use only functions from the tidyverse. Your code should work for strings
of this shape in general, not only for the rows below.

\begin{lstlisting}
notes <- tribble(
  ~ratio, ~words,
  "3/4",  "Undo",
  "5/2",  "into",
  "1/8",  "pretty",
  "9/9",  "ready",
  "2/5",  "inner",
  "7/10", "are.",
  "4/4",  "Lost"
)
\end{lstlisting}

\begin{enumerate}
\item (4 pts) Convert \texttt{ratio} from a character vector of fractions
to a numeric vector of decimal values. For example, \texttt{"3/4"} should
become \texttt{0.75}. No \texttt{stringr} functions and no regular expressions.

\item (4 pts) Use a single regular expression to find every word in
\texttt{words} that begins with \texttt{un} or \texttt{in}, ignoring case,
and replace that prefix with \texttt{re}. So \texttt{Undo} becomes
\texttt{redo}, and \texttt{into} becomes \texttt{reto}.

\item (2 pts) \textbf{Challenge.} Move the first character of each word to
the end, then append \texttt{ay}. Examples: \texttt{study} becomes
\texttt{tudysay}, and \texttt{aloe} becomes \texttt{loeaay}.
Build that version of each string in \texttt{words}.

For a little extra credit: keep an initial capital as a capital on the
new first letter (\texttt{John} becomes \texttt{Ohnjay}), and keep a
trailing period in place (\texttt{goodbye.} becomes \texttt{oodbyegay.}).
\end{enumerate}

\item \textbf{Tidying a score file} (20 pts).
The file \texttt{survey\_scores.csv} is on the course page under Data.
Each row is one person. The condition columns are
\texttt{respondent}, \texttt{city}, \texttt{year}, and \texttt{arm}.

\texttt{arm} is a string. It holds either one letter in parentheses,
such as \texttt{(A)}, or two letters in parentheses separated by a comma
and a space, such as \texttt{(A, B)}.

The other columns are scores. The name is a time, an underscore, and a
subject: \texttt{pre\_math}, \texttt{pre\_read}, \texttt{post\_math},
\texttt{post\_read}, \texttt{follow\_math}, and \texttt{follow\_read}.
The column \texttt{drop\_} is a score with no subject.

For each part, print the first six rows with \texttt{head()}.
The only non-tidyverse functions you may use in this question are
\texttt{head()} and \texttt{c()}.

\begin{enumerate}
\item (2 pts) Use \texttt{readr} to load the file, and view the first six rows.

\item (4 pts) Reshape the data so that, besides the four condition columns,
one column holds the score-column name and another column holds the score.

\item (3 pts) Split that name into a time column and a subject column.
For \texttt{drop\_}, the subject should be the empty string \texttt{""}.

\item (5 pts) Clean \texttt{arm}:
\begin{itemize}
  \item Remove every parenthesis.
  \item Replace the first letter with that same letter followed by an underscore.
  \item Remove every comma and the space after it.
\end{itemize}
For example, \texttt{(A, B)} becomes \texttt{A\_B}, and
\texttt{(A)} becomes \texttt{A\_}.

\item (3 pts) Using the cleaned \texttt{arm}, make two columns:
the first letter and the second letter. If \texttt{arm} had only one
letter, the second is \texttt{NA}. Keep both columns as character, not
numbers.

\item (3 pts) Make an appropriate plot of how time and subject relate to
the score. Leave out the \texttt{drop} rows, which have no subject.
\end{enumerate}
\end{enumerate}

\end{document}
